% Generated draft: requires mathematical and source review.
\documentclass[11pt]{article}
\usepackage[margin=25mm]{geometry}
\usepackage{fontspec}
\setmainfont{Latin Modern Roman}
\usepackage{amsmath,amssymb,mathtools,unicode-math}
\setmathfont{Latin Modern Math}
\usepackage{enumitem,needspace,xurl,hyperref}
\setlength{\parindent}{0pt}
\setlength{\parskip}{5pt}
\setlength{\emergencystretch}{3em}
\newcommand{\N}{\mathbb{N}}
\newcommand{\Z}{\mathbb{Z}}
\newcommand{\Q}{\mathbb{Q}}
\newcommand{\R}{\mathbb{R}}
\newcommand{\C}{\mathbb{C}}
\newcommand{\F}{\mathbb{F}}
\newcommand{\A}{\mathbb{A}}
\newcommand{\PP}{\mathbb{P}}
\newcommand{\E}{\mathbb{E}}
\newcommand{\eps}{\varepsilon}
\newcommand{\dd}{\,\mathrm{d}}
\begin{document}
\section*{UnsolvedMath ID: 30005394}
\textbf{Source problem code:} catalogue entry
\subsection{Definitions and mathematical statement}
Let \((G,X)\) be a Shimura datum of Hodge type with reflex field \(E\). Retain the introductory assumption that the adjoint group \(G^{\mathrm{ad}}\) is \(\Q\)-simple. Let \(p\) be a prime such that \(G_{\Q_p}\) is quasi-split and becomes split over an unramified extension of \(\Q_p\). Choose a hyperspecial subgroup \(U_p\subseteq G(\Q_p)\), a sufficiently small compact open subgroup \(U^p\subseteq G(\mathbb A_f^p)\), and a place \(v\mid p\) of \(E\). Here \(\mathbb A_f^p\) denotes the finite adeles away from \(p\), and \(U=U^pU_p\).

Let \(k=\overline{\kappa(v)}\), and write
\[
S=\operatorname{Sh}_{G,U}
\]
for the geometric special fibre over \(k\) of the canonical integral model over \(\mathcal O_{E,v}\) of the Shimura variety of level \(U\). The supplied excerpt attributes existence of this integral model to its reference [9], whose bibliographic entry is not supplied.

Let \(C\subseteq S\) be a central leaf in the sense of references [5] and [8] of the excerpt. In the usual Hodge-type interpretation, a central leaf fixes the isomorphism class of the associated \(p\)-divisible group together with its relevant additional tensor structure. The excerpt does not specify this structure or its precise isomorphism convention, so this description does not replace the missing referenced definition. Let \(h(G)\) denote the Coxeter number used in the excerpt; its convention for the reductive group \(G\) is not specified there. Stability under prime-to-\(p\) Hecke operators means preservation by all the corresponding prime-to-\(p\) Hecke correspondences, with the usual interpretation through changes of level.

\textbf{Statement explicitly supplied as a proved theorem (Theorem 1).} For every central leaf \(C\subseteq S\) and every nonempty reduced closed subvariety \(Z\subseteq C\) stable under the prime-to-\(p\) Hecke operators,
\[
p\geq h(G)\quad\Longrightarrow\quad Z=C.
\]
Closedness here is relative to \(C\). The assertion is recorded as a theorem in the supplied excerpt, not as an independently verified result.

\textbf{Unresolved intended problem.} The supplied dated curation describes the general Hodge-type Hecke orbit conjecture without the size restriction as remaining open. This suggests asking whether the displayed conclusion holds without \(p\geq h(G)\), while retaining the other hypotheses. However, the record does not supply an explicit conjecture establishing that this is its exact intended formulation. Its statement-verification field marks the problem statement as unrecoverable because the extracted passage states a proved theorem. Consequently, removal of the bound is only a suggested target, not a recovered source conjecture.

The excerpt calls an earlier, omitted conjecture the Hecke orbit conjecture for \(\mathcal A_g\), the Siegel moduli space, and attributes its discrete part to Theorem C of [10], with related results in [7]. It also attributes the Siegel case of the displayed theorem, without any condition on \(p\), to Chai--Oort: the continuous part to their then-forthcoming book [4] and the discrete part to [3]. These attributions are supplied claims. The excerpt says their proofs depend on hypersymmetric points in Newton strata and therefore do not generalize directly, since such points usually do not exist in the more general setting; the passage ends mid-sentence.
\subsection{Short English statement}
The supplied theorem says that, when \(p\geq h(G)\), every nonempty reduced closed subset of a central leaf that is stable under all prime-to-\(p\) Hecke correspondences equals the whole leaf. The curation suggests removing this bound in the Hodge-type setting, but the exact open conjecture is missing. The record labels the topic partially solved and the extracted problem statement unrecoverable; this is a statement-only draft requiring mathematical review.
\subsection{Sources}
\begin{itemize}
\item Local registry: \texttt{data/unsolvedmath/problems.json}, with \texttt{catalog/problems.json} as the supplementary registry; supplied source record for permanent UnsolvedMath ID \(30005394\), titled ``Central Leaves Under Prime-to-\(p\) Hecke Correspondences,'' exact source code \texttt{catalogue entry}. No independent inspection of these registry files is asserted.
\item Supplied report citation: \emph{Arithmetic of Shimura Varieties} (2023), Oberwolfach Reports, DOI 10.4171/owr/2023/5, pp.\ 261--326. \url{https://doi.org/10.4171/owr/2023/5}. The record describes Oberwolfach Reports as workshop reports documenting mathematical discussions and open problems; no individual proposer is supplied.
\item Supplied literature source: DAddezio and van Hoften, \emph{Hecke orbits on Hodge-type Shimura varieties}. \url{https://arxiv.org/abs/2205.10344}.
\item Supplied literature source: van Hoften, \emph{Ordinary Hecke orbit conjecture}. \url{https://arxiv.org/abs/2112.12422}.
\item Supplied literature source: \emph{Oberwolfach Report 5/2023}. \url{https://doi.org/10.4171/owr/2023/5}.
\item Literature assessment dated \texttt{2026-08-22}: the supplied curation labels the topic partially solved and states that the discrete part, the continuous part under a Coxeter-number bound, and further ordinary and Siegel cases are established, while the general Hodge-type conjecture without the size restriction remains open. This is supplied dated curation derived from the OWR record, not an independently verified literature assessment; its time-sensitive status and cited claims require independent verification.
\item The excerpt attributes the unrestricted Siegel result to Chai--Oort, with its continuous part in the then-forthcoming book [4] and discrete part in [3]. It also cites Theorem C of [10], related results in [7], central-leaf constructions in [5] and [8], and integral-model existence in [9]. The bibliography identifying these numbered references is absent from the supplied record.
\end{itemize}
\end{document}
